\chapter{Dataset and Tools}
\label{chap4}

This chapter describes the data and the software used to build the system. It first presents the CUAD contract dataset that the system retrieves over and from which the evaluation questions are drawn, including its format and the properties that affect how the results should be interpreted. It then introduces the main libraries and frameworks used in the implementation, from the retrieval and generation stack through to the observability and evaluation tools.

\section{The CUAD Dataset}

CUAD (Contract Understanding Atticus Dataset) was introduced by Hendrycks et al.\ \cite{hendrycks2021cuad} as a large-scale benchmark for machine learning on legal contracts. It contains 510 commercial contracts collected from SEC-EDGAR, the public filings database of the U.S.\ Securities and Exchange Commission. Expert lawyers labelled each contract against 41 clause categories that practitioners commonly flag during contract review, including Governing Law, Non-Compete, IP Ownership Assignment, Expiration Date, and Termination for Convenience, among others. The dataset contains more than 13{,}000 manually created annotations in total. Of the 41 categories, 33 require a yes/no answer indicating whether a given clause exists in the contract, while the remaining 8 require an extractive answer such as a state name, a date, or an entity name. This distinction matters for evaluation: the system must handle both the detection of a clause and the extraction of its exact wording.
\subsection{Data Format}

CUAD follows the SQuAD extractive reading comprehension format \cite{rajpurkar2016squad}. Each record pairs a question with a passage from one contract, and the expected answer is a verbatim span extracted from that passage. A boolean field \texttt{is\_impossible} marks whether the queried clause is absent from the contract, enabling evaluation of the system's ability to recognise unanswerable questions. Alongside the annotated JSON file, the dataset distributes all 510 contract texts as plain text files.

\subsection{Usage in This Thesis}
\label{sec:question-sampling}
The 510 contract texts were loaded, split into chunks, and stored in three separate ChromaDB vector databases, one per chunking strategy. The chunking strategies themselves are described in
Chapter~\ref{chap5}.

For evaluation, a stratified random sample of questions was drawn from the dataset
annotations. The sample consisted of 20 questions with a fixed random seed to ensure
reproducibility: 70\,\% were answerable questions (\texttt{is\_impossible=False}) and the remaining
30\,\% were unanswerable (\texttt{is\_impossible=True}). This split ensures that the
evaluation covers both the system's ability to extract the relevant clause text and its
ability to correctly decline when no matching clause exists.

Before each question is submitted to the pipeline, the title of the corresponding
contract is prepended. This allows the retrieval step to focus on the correct
document rather than searching across all 510 contracts simultaneously.

\subsection{Dataset Limitations and Evaluation Considerations}

Working with CUAD in an automated evaluation context revealed several properties of
the dataset that affect how results should be interpreted.

\textbf{Non-conversational question format.}
Every question in CUAD follows a fixed boilerplate template: ``Highlight the
parts (if any) of this contract related to `[Category]' that should be reviewed by a
lawyer. Details: [description].'' This is a legal review instruction, not a natural
question. When RAGAS computes Answer Relevancy, it generates reverse questions from
the system's answer and compares them to the original question using embeddings.
Because the language model answers as if the question were a factual query, the
reverse questions end up semantically distant from the CUAD boilerplate, and the
metric penalises the answer even when it is factually correct. This reflects a
structural mismatch between the CUAD format and the RAGAS evaluation methodology,
not a failure of the system, and it should be kept in mind when interpreting Answer
Relevancy scores.

\textbf{Semantic gap between category names and clause text.}
CUAD category names are short, standardised labels (\textit{Expiration Date},
\textit{Minimum Commitment}, \textit{Governing Law}), whereas the corresponding text
in contracts uses varied legal phrasing: a contract may express expiration as
\textit{``shall terminate upon 90 days written notice''}, or state the governing law
as \textit{``construed in accordance with the laws of the State of Delaware''}. These
pairs are legally equivalent but semantically distant in embedding space, so a
semantic-only retriever can miss the relevant clause even when it is present.This motivated the addition of hybrid retrieval, in which a BM25 keyword index runs alongside the dense retriever and its queries are expanded with legal synonyms for each clause category. Both are described in detail in Chapter~\ref{chap5}.


\textbf{Structured data and annotation quirks.}
Some ground truth answers in CUAD are formatted as tables containing payment
schedules with redacted amounts. Tabular content with few meaningful tokens embeds
poorly and is difficult for any semantic retriever to rank highly. Additionally,
CUAD annotators occasionally marked a termination clause as the ground truth for an
\textit{Expiration Date} question in contracts that carry no fixed expiration date.
In such cases the question and its own ground truth are semantically misaligned,
which can produce misleading Context Recall scores. Both issues are properties of
the dataset rather than the retrieval system.


\section{Implementation Tools}

This section introduces the main libraries and frameworks used to implement the system, grouped by their role in the pipeline: orchestration, storage and retrieval, generation, the user interface, observability, and evaluation.

\subsection{LangChain}

LangChain \cite{chase2022langchain} is a Python framework for building applications
around language models. It provides abstractions for document loading, text
splitting, vector stores, prompt templates, and retrieval chains. In this project,
LangChain handles the entire pipeline: loading the contract files, splitting them
using three different strategies, interfacing with the ChromaDB vector stores,
constructing the prompt templates for question answering and conversation
condensation, and connecting all components through its expression language.

\subsection{ChromaDB}

ChromaDB \cite{chroma2023} is an open-source vector database designed for local
deployment. It stores document embeddings in an HNSW index \cite{malkov2018hnsw},
which supports approximate nearest-neighbour search in sub-linear time. Embeddings
and metadata are persisted to disk between runs. Three separate databases were
created for this project, one for each chunking strategy, so that each experimental
configuration queries its own index and results across configurations remain
directly comparable.

\subsection{Hugging Face Transformers and bitsandbytes}

The Hugging Face Transformers library \cite{wolf2020transformers} was used to load
Llama~3.1~8B Instruct and Mistral~7B Instruct locally from the Hugging Face Hub.
Both models were loaded with 4-bit NF4 quantization via bitsandbytes
\cite{dettmers2022qlora}, reducing memory consumption from approximately 16\,GB in half precision
to roughly 5\,GB per model with minimal quality loss on extractive tasks. Model
weights are distributed across the available GPU automatically. Each model is
subsequently wrapped so that the tokenizer's native chat template is applied
correctly, ensuring the instruction-tuned model receives the expected input format
without echoing the prompt back into its output.


\subsection{Sentence Transformers}

All document chunks and queries were embedded using the all-MiniLM-L6-v2 model from
the Sentence Transformers library \cite{reimers2019sbert}. The model produces
384-dimensional vectors optimised for semantic similarity tasks including passage
retrieval. It was selected for its low computational cost relative to larger
alternatives while still providing reliable retrieval quality for English legal text.

\subsection{BM25 Keyword Retrieval}

In addition to semantic search, the system uses BM25 \cite{robertson2009bm25} for keyword-based retrieval. As noted earlier in this chapter, semantic search can fail when the query and the contract text use different but legally equivalent phrasing. BM25 bridges this gap by matching exact and near-exact terms regardless of semantic distance. The BM25 index is built from the documents stored in the active vector database, and its results are merged with the semantic search hits before being passed to the LLM.

\subsection{Langfuse}

Langfuse \cite{langfuse2023} is an observability platform for LLM applications.
Each call to the RAG pipeline creates a trace that records the user's question, the
reformulated retrieval query, the retrieved chunks, the generated answer, and the
latency of each processing step. In this thesis, Langfuse Cloud is used during the
experimental evaluation to collect per-configuration metrics such as generation
latency and token consumption across all six pipeline configurations, enabling a
direct comparison between the two language models and the three chunking strategies.


\subsection{Streamlit}

Streamlit \cite{streamlit2019} is a Python library for building interactive web
applications. It was used to develop the main user-facing interface of the system:
a browser-based chat application where a user can select a chunking strategy, a
language model, and a memory type, and then hold a natural-language conversation
with the RAG pipeline about the loaded contracts. This interface is the primary
deliverable of the thesis, providing direct access to the system without requiring
any programming knowledge. Language models and vector stores are cached so that
switching one configuration option does not trigger a full reload of the pipeline.


\subsection{Docker}

The full application is containerised with Docker. A dedicated configuration file
installs PyTorch with CUDA support first, then all remaining dependencies. The web
interface is exposed on a fixed port. Langfuse Cloud is accessed through API keys
supplied as environment variables, so no local observability service runs alongside
the container. Containerisation ensures the environment is fully reproducible across different
machines.

\subsection{RAGAS}

RAGAS \cite{es2024ragas} is an evaluation library built specifically for RAG
systems. It computes Faithfulness, Answer Relevancy, Context Precision, and
Context Recall automatically using a judge language model, without requiring manual
annotation of each response. The evaluation pipeline collects the system's answers
and the retrieved context, passes them to the RAGAS evaluator, and writes the
per-question scores together with their averages to a timestamped CSV file for
further analysis.

